
\documentclass{article}
\usepackage{tikz}
\usepackage{verbatim}

% \usepackage{fontspec}
 
% \setmainfont{Times New Roman}


% drawing parameters
\newcommand{\circdist}{1.2}  % distance from origin to center of circles
\newcommand{\circrad}{7/4} % radius of the circles
\newcommand{\circlethickness}{6mm} % uh, thickness of the circles

% distance from the origin to the three "interior" intersections
\pgfmathsetmacro{\intrad}{sqrt((\circrad)^2 - 3*(\circdist)^2/4) - \circdist/2}

% distance from the origin to the three "exterior" intersections
\pgfmathsetmacro{\extrad}{sqrt((\circrad)^2 - 3*(\circdist)^2/4) + \circdist/2}


\definecolor{FFgelb}{HTML}{FFB400}
\definecolor{FFmagenta}{HTML}{DC0067}
\definecolor{FFblau}{HTML}{009EE0}

% so we can just specify an angle and get the correct color for the circle
\colorlet{180}{FFmagenta}
\colorlet{40}{FFmagenta}
\colorlet{20}{FFmagenta}
\colorlet{300}{FFmagenta}
\colorlet{180}{FFmagenta}

% draws one of our circles
\newcommand{\mycircle}[1]{%
  \draw[FFmagenta, double distance=\circlethickness, double=#1]
  (#1:\circdist) circle (\circrad);}
  
 % draws one of our circles
\newcommand{\freifunks}[1]{%
  \draw[FFmagenta, double distance=2 mm, double=#1]
  (#1:\circdist) circle (7/5);}
  
\newcommand{\freifunkleft}[1]{%
  \draw[FFmagenta, double distance=2 mm, double=#1]
  (#1:\circdist) 
  circle (7/5);}

\begin{document}

\begin{center}

  % \begin{tikzpicture}
  %   \draw[yellow,line width=1cm] (0,0) circle[radius=2cm];
  % \end{tikzpicture}

\begin{tikzpicture}[scale=2]

  % draw the circles
%  \foreach \angle in {180,20}
%  {
%    \mycircle{\angle}
%  }
	\freifunkleft{180}
	\mycircle{20}
  \freifunks{20}

  \filldraw[fill=FFgelb, color=FFgelb]
(0,2) -- (3,2) -- (2.5,1) -- (3.5,1) -- (4.5,2.5) %
-- (3.5,4) -- (2.5,4) -- (3,3) --(0,3)
-- (0,2) ;

  % \draw[black](0,0) {some text}

  % we need to re-do two of the intersections so that the rings
  % interlock and aren't just piled on top of one another
  \foreach \angle/\rad in {60/\intrad, 240/\extrad}
  {
    \begin{scope}
      \clip (\angle:\rad) circle (5/4*\circlethickness);
      %                          you may need to adjust the 5/4 above so
      %                          the clipping area is big enough
     % \mycircle{300}
     % \mycircle{180}
    \end{scope}
  }
\end{tikzpicture}


  \begin{tikzpicture}[overlay,remember picture]
    \node [
      fill=FFmagenta,% Farbe des Randstreifens
      text=white,% Textfarbe
      font=\normalfont\bfseries,% Einstellungen für die Schrift
      inner xsep=1em, % Abstand des Textes von unten
      % maximale Textbreite = Papierhöhe - 2*Abstand des Textes von unten:
      text width={\dimexpr\paperheight-2em\relax},
      minimum height=17mm,% Breite des Randstreifens
      anchor=north west,
      rotate=90
      ]
      at (current page.south west)
      {fulda.freifunk.net};
  \end{tikzpicture}%
    \begin{tikzpicture}[overlay,remember picture]
    \node [
      fill=FFblau,% Farbe des Randstreifens
      text=white,% Textfarbe
      font=\normalfont\bfseries,% Einstellungen für die Schrift
      inner xsep=1em, % Abstand des Textes von unten
      % maximale Textbreite = Papierhöhe - 2*Abstand des Textes von unten:
      text width={\dimexpr\paperheight-2em\relax},
      minimum height=17mm,% Breite des Randstreifens
      anchor=south east,
      rotate=90
      ]
      at (current page.north east)
      {fulda.freifunk.net};
  \end{tikzpicture}%



fulda.freifunk.net


\begin{tikzpicture}
\fill[green,opacity=0.5] (4.25,0) circle (0.25);
\end{tikzpicture}





\begin{tikzpicture}
\filldraw[fill=FFgelb, color=FFgelb]
(0,2) -- (3,2) -- (2.5,1) -- (3.5,1) -- (4.5,2.5) %
-- (3.5,4) -- (2.5,4) -- (3,3) --(0,3)
-- (0,2) ;

%----------------------

% \draw[] (0,3)--(-1,1)--(2,1)--(2,2)--(4,0) node[xshift=-7ex, align=center]{gewünschte \\ Form};
%\draw[yscale=-1] (0,0)--(-1,1)--(2,1)--(2,2)--(4,0);

\end{tikzpicture}








\end{center}

\end{document}